\documentclass[11pt]{article}
\usepackage{amsmath,amssymb,amsthm,geometry,hyperref}
\geometry{margin=0.85in}
\newtheorem{proposition}{Proposition}
\title{A counterexample to Conjecture 00000002618}
\author{gaochengzhecpu (AI-assisted submission)}
\date{October 3, 2026}
\setlength{\emergencystretch}{3em}
\begin{document}
\maketitle

The conjecture states, in particular, that a finite lattice is
semidistributive if and only if each element has a unique canonical join
representation. We disprove the reverse implication. All definitions below
are the standard lattice-theoretic ones. In particular, semidistributivity
requires both the join and the meet semidistributive laws.

\paragraph{Definitions.}
For a subset $A$ of a finite lattice $L$, a join representation of $w$ is
an equality $w=\bigvee A$, with $\bigvee\varnothing=0$.
It is irredundant when no proper subset of $A$ has join $w$.
A subset $A$ \emph{join-refines} $B$ if for every $a\in A$ there exists
$b\in B$ with $a\leq b$.
A canonical join representation is an irredundant join representation
that join-refines every irredundant join representation of the same
element. This is the standard definition used, for example, by Barnard
\cite{Barnard}. The argument below verifies the stronger property of
refining every join representation.

The join and meet semidistributive laws, respectively, are
\[
 x\vee y=x\vee z\ \Longrightarrow\
 x\vee(y\wedge z)=x\vee y,\qquad
 x\wedge y=x\wedge z\ \Longrightarrow\
 x\wedge(y\vee z)=x\wedge y.
\]
A semidistributive lattice satisfies both.

\paragraph{The lattice.}
Let $E=\{a,b,c\}$ and let
\[
 L=\{\varnothing,\{a\},\{b\},\{a,b\},\{c\},\{a,c\},E\}
\]
be ordered by inclusion. Thus the subset $\{b,c\}$ is omitted from the
Boolean lattice of subsets of $E$.
This family is closed under intersections, contains $\varnothing$ and
$E$, and hence is a lattice. Its meet is intersection. Its join is the
union unless that union is $\{b,c\}$, in which case the join is $E$.
For clarity write its elements as $0,A,B,D,C,F,1$, in the order displayed
above, so $D=A\vee B$, $F=A\vee C$, and $1=B\vee C$.

\begin{proposition}
Every element of $L$ has a unique canonical join representation.
\end{proposition}
\begin{proof}
The proposed representations are
\[
\begin{array}{c|ccccccc}
 w&0&A&B&D&C&F&1\\ \hline
 \operatorname{Can}(w)&\varnothing&\{A\}&\{B\}&\{A,B\}
 &\{C\}&\{A,C\}&\{B,C\}.
\end{array}
\]
Each has join $w$ and is irredundant: in the two-element cases either
single remaining atom is strictly below $w$; in the one-element cases
removing the atom leaves join $0$. The empty representation is
irredundant.

We verify refinement against an arbitrary join representation $R$ of $w$.
For $w=0$ refinement by the empty set is vacuous.
If $w$ is an atom, every element of $R$ is below $w$, and $R$ must contain
$w$, since otherwise its join would be $0$.
For $w=D$, if no element of $R$ lies above $A$, then every element of $R$
belongs to $\{0,B\}$, whence $\bigvee R\leq B<D$, a contradiction.
The same reasoning with $A$ and $B$ exchanged gives an element above $B$.
Thus $\{A,B\}$ refines $R$. The argument for $w=F$ is identical using
$A,C$.
For $w=1$, if no element of $R$ lies above $B$, then every element of $R$
is in $\{0,A,C,F\}$, all below $F<1$, a contradiction.
If none lies above $C$, all elements of $R$ are below $D<1$, again a
contradiction. Therefore $\{B,C\}$ refines $R$.

Finally canonical join representations are unique as subsets.
Indeed any irredundant representation is an antichain: if distinct
$x\leq y$ both occurred, removing $x$ would leave the same join.
If two canonical representations $S,T$ refine one another, then for
$x\in S$ there exist $y\in T$ and $z\in S$ with $x\leq y\leq z$.
The antichain property of $S$ gives $x=z$, and antisymmetry gives
$x=y$, so $x\in T$. By symmetry $S=T$.
\end{proof}

\paragraph{Failure of semidistributivity.}
Take $x=A$, $y=B$, $z=C$. Then
\[
 A\wedge B=0=A\wedge C,\qquad
 A\wedge(B\vee C)=A\wedge1=A\ne0.
\]
The meet semidistributive law fails. This proves that the lattice is not
semidistributive, despite the unique canonical join representations of
all its elements, and therefore disproves the stated equivalence.
The further congruence-separation assertion in the source cannot repair
this already false conjunct. No interpretation of that additional
assertion is needed.

\paragraph{Scope of the counterexample.}
The familiar correct finite-lattice criterion characterizes
\emph{join-semidistributivity} by existence of canonical join
representations. Full semidistributivity additionally requires the dual
meet condition; this distinction is explicit in \cite{Barnard}.
Our lattice in fact satisfies the join semidistributive law, as is also
checked by the formal proof. We use the standard meaning of the
unqualified word \emph{semidistributive}, consistent in the source's
English and Chinese statements.

\paragraph{Formal verification and semantic correspondence.}
The accompanying file \texttt{lean/Main.lean} uses Lean 4.19.0 and
\texttt{Std}, with no external mathematical library.
\texttt{LatticeData} contains a partial order, a bottom element, and
both universal properties for binary joins and meets. The concrete
object \texttt{seven} verifies every one of these axioms on
\texttt{Fin 7}. Its labels encode the displayed subsets by their
three-bit membership masks, with label $6$ representing $E$ rather than
the missing subset $\{b,c\}$.

\texttt{JoinRep} asserts the least-upper-bound property for an
arbitrary predicate $S:L\to\mathrm{Prop}$, so it is exactly
$w=\bigvee S$. \texttt{Irredundant}, \texttt{Refines}, and
\texttt{CanonicalJoin} implement the definitions above, quantifying over
arbitrary subsets.
All $2^7=128$ subsets have finite binary codes.
Crucially, \texttt{every\_subset\_encoded} proves that every
predicate-valued subset equals a decoded code.
Thus finite verification does not restrict the mathematical
representations to a selected family.
\texttt{irredCode\_iff} transfers the finite irredundancy check to its
unrestricted definition.
The kernel checks all seven representation certificates; the general
antichain argument \texttt{canonical\_unique} then supplies uniqueness.

\texttt{counterexample} states that this genuine lattice has unique
canonical join representations for every element and is not
semidistributive.
\texttt{ClaimedCriterion} is the necessary reverse implication of the
source for all finite labelled lattices, and
\texttt{conjecture2618\_false} proves its negation by specializing to
\texttt{seven}.
All finite checks use kernel-reduced \texttt{decide}; there is no
\texttt{native\_decide}, admitted result, or additional axiom.

\begin{thebibliography}{1}
\bibitem{Barnard}
E. Barnard, \emph{The canonical join complex},
Electronic Journal of Combinatorics \textbf{26}(1) (2019), P1.24.
\url{https://www.combinatorics.org/ojs/index.php/eljc/article/download/v26i1p24/pdf/}
\end{thebibliography}
\end{document}
